\chapter{Implementace systému NPC}
\label{sec:implementace}

Implementace byla rozdělena do několika navazujících částí. 
Nejprve je popsána realizace základní třídy NPC a jejích odvozených variant, poté implementace AI Controlleru, Blackboardu a behaviorálního stromu. 
Následující části se věnují konkrétním Taskům a Services, správě hlídkových tras, zpracování percepčních podnětů a komunikaci mezi NPC. 
Závěr kapitoly shrnuje vybraná implementační rozhodnutí a problémy, které bylo nutné při vývoji řešit. 

\section{Implementace BP\_NPC}
\label{sec:implementace-bp-npc}

Třída \texttt{BP\_NPC} je odvozena od třídy \texttt{BP\_ThirdPersonCharacter} a představuje základní implementaci všech nehráčských postav ve hře. 
Jejím hlavním účelem je realizace konkrétních akcí ve světě, zatímco rozhodovací logika je delegována na AI Controller a behaviorální strom. 
Třída tak slouží především jako vykonavatel příkazů generovaných vyššími vrstvami systému.

V rámci třídy jsou implementovány základní mechaniky pohybu, animací, útoku a interakcí s okolím. 
Aktualizace stavových proměnných souvisejících s pohybem probíhá periodicky v rámci události \texttt{Tick}, kde je z komponenty \texttt{CharacterMovement} získána aktuální rychlost postavy. 
Její velikost je následně ukládána do proměnné \texttt{Speed}, která slouží jako vstup pro animační systém. 
Současně je pomocí funkce \texttt{IsFalling} určováno, zda se postava nachází ve vzduchu, a výsledek je ukládán do proměnné \texttt{IsJumping}. 
Tyto proměnné umožňují animačnímu blueprintu plynule přepínat mezi jednotlivými animačními stavy.

Útočné chování je realizováno prostřednictvím vlastní události \texttt{Attack}, která je volána z behaviorálního stromu, k čemuž slouží \textit{task} \texttt{BTT\_Attack}. 
Po jejím spuštění je nastavena proměnná \texttt{IsAttacking} a přehrána odpovídající animační montáž pomocí uzlu \texttt{Play Montage}. 
Po dokončení nebo přerušení animace je proměnná \texttt{IsAttacking} opět resetována. 
Za povšimnutí stojí systém notifikací. V rámci animace (respektivě animační montáže) může vývojář vložit události na konkrétní snímky 
animace, čímž je mu dovoleno velmi flexibilně ovlivňovat dění v herním světě v průběhu animace.
Pro využití systému notifikací stačí vložit notifikaci na časové ose animace a zvolit buď jeden z předpřipravených 
událostí (např. \textit{Play Sound}), nebo vlastní událost, jak jde vidět na obrázku~\ref{fig:notify_timeline}.
Obrázek \ref{fig:notify_blueprint} poté ukazuje, jak taková událost může být implementována.
Tato konkrétní událost při zásahu hráče teleportuje hráče zpátky na \texttt{PlayerStart} (začátek levelu).

\begin{figure}[H]
    \centering

    \includegraphics[width=\textwidth]{Figures/AN_OnHitEvent.png}
    \caption{Umístění události \texttt{AN\_HitEvent} v animační montáži}
    \label{fig:notify_timeline}
\end{figure}

\begin{figure}[H]
    \centering

    \includegraphics[width=\textwidth]{Figures/AN_OnHitEventBlueprint.png}
    \caption{Implementace události \texttt{AN\_HitEvent} v blueprintu}
    \label{fig:notify_blueprint}
\end{figure}

Manipulace se zbraní je řešena pomocí dvojice událostí \texttt{EquipWeapon} a \texttt{UnequipWeapon}. 
Událost \texttt{EquipWeapon} přehrává animační montáž tasení zbraně a následně vytváří instanci zbraně pomocí \texttt{SpawnActor}. 
Vytvořený objekt je připojen ke kostře postavy na konkrétní socket. 
Reference na vytvořenou zbraň je uložena do proměnné \texttt{WeaponActor} a příznak \texttt{IsWieldingWeapon} indikuje, že postava aktuálně drží zbraň. 
Opačný proces je implementován v události \texttt{UnequipWeapon}, kde je zbraň odpojena a odstraněna.

Interakce s okolními objekty je realizována prostřednictvím kolizního systému. 
Při vstupu do kolize s objektem implementujícím rozhraní \texttt{BPI\_Interactible} je tato skutečnost detekována v události \texttt{OnComponentBeginOverlap}. 
Následně je získán AI Controller dané postavy a prostřednictvím jeho rozhraní je aktualizována odpovídající hodnota v Blackboardu, čímž je behaviorálnímu stromu signalizována možnost interakce. 
Při ukončení kolize (\texttt{OnComponentEndOverlap}) je tato informace opět odstraněna.

Třída dále obsahuje pomocnou funkci \texttt{SetMovementSpeed}, která na základě hodnoty výčtového typu nastavuje maximální rychlost pohybu postavy. 

Konkrétní postavy \texttt{BP\_Grandpa} a \texttt{BP\_Granny} dědí z třídy \texttt{BP\_NPC} a upravují především vizuální reprezentaci a animační parametry. 
Sdílená logika chování zůstává v třídě rodiče, čímž je minimalizována duplicita a zjednodušena případná rozšiřitelnost systému.

\section{Implementace AI Controlleru (IAC\_NPC)}
\label{sec:implementace-iac-npc}

Třída \texttt{IAC\_NPC}, odvozená od~\texttt{AIController}, představuje hlavní řídicí komponentu NPC. 
Jejím úkolem je zprostředkovat komunikaci mezi percepčním systémem, behaviorálním stromem a herním světem, a zároveň udržovat aktuální stav 
postavy prostřednictvím \textit{blackboardu}.

\subsection{Inicializace a napojení na behaviorální strom}

Po převzetí řízení nad postavou (událost \texttt{OnPossess}) dochází k inicializaci celého systému.
Nejprve jsou vyhledány všechny instance tříd reprezentujících hlídkové trasy a uloženy do interní kolekce.
Tento krok umožňuje dynamicky reagovat na změny ve scéně bez nutnosti pevného navázání tras v editoru.

Následně je spuštěn behaviorální strom \texttt{BT\_NPC}, přičemž před jeho spuštěním jsou do \textit{blackboardu}
zapsány výchozí hodnoty. Jedná se zejména o:
\begin{itemize}
    \item referenci na řízenou postavu (\texttt{SelfActor}),
    \item pravděpodobnost změny hlídkové trasy,
    \item pravděpodobnost interakce s objektem,
    \item počáteční hlídková trasa.
\end{itemize}

Inicializační fáze tak zajišťuje, že behaviorální strom pracuje nad konzistentní sadou dat a může okamžitě začít rozhodovací proces.

\subsection{Zpracování percepce}

Pro vytvoření dojmu smyslového vnímání u NPC se využívá komponenta \texttt{AIPerception}, která agreguje informace o okolních aktérech
a vyvolává událost \texttt{OnTargetPerceptionUpdated} při každé změně percepčního stavu.

Zpracování této události probíhá tak, že pro každý aktualizovaný aktér jsou získána
percepční data, jednotlivé podněty jsou analyzovány podle typu smyslu
a na základě výsledku jsou následně aktualizovány klíče v \textit{blackboardu}.

Pro rozlišení typu podnětu je využita pomocná metoda \texttt{CanSenseActor}, jejíž implementace 
je zobrazena na obrázku~\ref{fig:can_sense_actor}.
Metoda nejprve získá percepční data pro daného aktéra pomocí uzlu \texttt{Get Actors Perception}. 
Tato data jsou následně rozložena na jednotlivé stimuly, ze kterých se percepční informace skládá.

Následně je ověřeno, zda některý ze stimulů odpovídá požadovanému typu smyslu.

Pokud dojde ke shodě typů, metoda vrací informaci o úspěšném zaznamenání podnětu spolu s odpovídajícím stimulem. 
V opačném případě pokračuje vyhodnocování dalších stimulů, dokud není nalezen odpovídající podnět
nebo není seznam vyčerpán.

Pro obě NPC je nastaven zrak (\texttt{AISense\_Sight}) jako dominantní smysl.
Dominance vjemu znamená, že pokud je v rámci jedné aktualizace percepce
zaznamenáno více podnětů současně, je prioritně vyhodnocován právě tento smysl.

\begin{figure}[H]
    \centering

    \includegraphics[width=\textwidth]{Figures/CanSenseActor.png}
    \caption{Implementace metody CanSenseActor}
    \label{fig:can_sense_actor}
\end{figure}


\subsubsection{Zrakové podněty}

Zrakové podněty jsou zpracovávány funkcí \texttt{HandleSensedSight}, která navazuje na funkci \texttt{CanSenseActor}. 
Tato funkce přijímá referenci na detekovaného aktéra a hodnotu určující, zda byl aktér aktuálně úspěšně zaznamenán. 
Na rozdíl od sluchového podnětu zde nestačí pouze uložit polohu stimulu, protože zrak může detekovat různé typy objektů s odlišným významem pro chování NPC.

Prvním krokem je ověření, zda je detekovaným aktérem hráč. 
Pokud ano a podnět je úspěšně zaznamenán, funkce zjistí aktuální stav NPC pomocí \texttt{GetCurrentState}. 
Nenachází-li se NPC ve stavu \texttt{Attacking}, je zavolána funkce \texttt{SetStateAsAttacking}. 
Ta zapíše referenci na hráče do klíče \texttt{AttackTarget} a nastaví klíč \texttt{State} na hodnotu \texttt{Attacking}. 
Tím je behaviorálnímu stromu předána informace, že má přejít do větve útočného chování.

Druhou částí zpracování zrakového podnětu je práce s interaktivními objekty. 
Funkce ověřuje, zda detekovaný aktér implementuje rozhraní \texttt{BPI\_Interactible}. 
Pokud ano, je porovnán s hodnotou aktuálně uloženou v klíči \texttt{SeenInteractible}. 
V případě platného nově zaznamenaného objektu je reference na tento objekt uložena do \textit{blackboardu}. 
Pokud naopak NPC daný objekt již nezaznamenává a tento objekt odpovídá aktuálně uložené hodnotě, je klíč \texttt{SeenInteractible} vymazán.

Tato logika umožňuje, aby jeden percepční mechanismus sloužil jak pro reakci na hráče, tak pro průběžné sledování objektů, se kterými může NPC později interagovat. 
Rozhodnutí, zda s objektem skutečně interagovat, však není provedeno přímo v AI Controlleru, ale až v příslušné části behaviorálního stromu.


\subsubsection{Sluchové podněty}

Sluchové podněty jsou zpracovávány funkcí \texttt{HandleSensedSound}. 
Oproti zrakovým podnětům je tato část implementace jednodušší, protože zvukový podnět je reprezentován pouze polohou, nikoliv konkrétním cílovým aktérem.

Funkce nejprve zjistí aktuální stav NPC pomocí \texttt{GetCurrentState}. 
Pokud se postava nenachází ve stavu \texttt{Investigating}, je zavolána funkce \texttt{SetStateAsInvestigating}. 
Ta nastaví klíč \texttt{State} na hodnotu \texttt{Investigating} a uloží polohu zvukového podnětu do klíče \texttt{PointOfInterest}. 
Behaviorální strom tak při dalším vyhodnocení přejde do větve vyšetřování a NPC se přesune k místu, odkud byl podnět zaznamenán.

Pokud se NPC nachází ve stavu \texttt{Attacking}, ke změně na vyšetřovací stav nedochází. 
Tím je zajištěno, že zvukový podnět nepřeruší probíhající útočné chování.

\subsection{Správa stavů}

Aktuální stav NPC je reprezentován výčtovým typem \texttt{E\_AIState} a uložen v \textit{blackboardu}. 

Pro práci se stavem jsou implementovány pomocné metody:
\begin{itemize}
    \item \texttt{GetCurrentState} – vrací aktuální stav,
    \item \texttt{SetStateAsPassive} – nastaví pasivní chování,
    \item \texttt{SetStateAsAttacking} – nastaví útočný režim,
    \item \texttt{SetStateAsInvestigating} – nastaví režim vyšetřování.
\end{itemize}

Každá z těchto metod provádí nejen změnu stavu, ale i odpovídající aktualizaci dalších klíčů (např. cíle útoku nebo bodu zájmu).
Tím je zajištěno, že behaviorální strom vždy pracuje s konzistentními daty.

\subsection{Interakce s objekty}

Součástí implementace je i podpora interakcí s objekty ve scéně.
Při detekci aktéra, který implementuje rozhraní pro interakci, je jeho reference uložena do klíče \texttt{CurrentInteractionTarget}.
Pokud se jedná o nový objekt (odlišný od posledního), je tato změna propagována do \textit{blackboardu}.

Další metody umožňují reagovat na konkrétní události:
\begin{itemize}
    \item \texttt{NotifyItemPickUp} – zaznamená typ drženého objektu,
    \item \texttt{NotifyItemPlaced} – vymaže informaci o drženém objektu.
\end{itemize}

Tyto informace jsou následně využívány behaviorálním stromem pro rozhodování o dalším postupu.

\section{Implementace stromu chování (BT\_NPC)}
\label{sec:implementace-bt-npc}

Strom chování \texttt{BT\_NPC} se větví do tří hlavních větví (\texttt{Attacking}, \texttt{Investigating}, \texttt{Passive}).

Chování NPC je dále ovlivněno stavem několika klíčů v \textit{blackboard}, které reprezentují aktuální fázi hry.
Na základě těchto hodnot dochází ke změnám v dostupných hlídkových trasách.
\begin{itemize}
    \item \texttt{Fáze hledání} -- hráč hledá klíč od venkovní boudy,
    \item \texttt{Fáze pátrání} -- hráč odemknul venkovní boudu,
    \item \texttt{Fáze útěku} –- hráč se snaží utéct z pozemku.
\end{itemize}

Přechody mezi jednotlivými fázemi jsou řízeny pomocí \textit{blackboard klíčů} \texttt{SearchShed} a \texttt{IsShedOpen}, které odrážejí postup 
hráče ve hře.

Tato struktura přirozeně vyjadřuje prioritizaci chování: útok má přednost před~pátráním a~pátrání má přednost před~hlídkováním.
NPC zmenšují počet hlídkovacích tras, které si vybírají k hlídkování a zvyšují svůj pobyt venku, aby zamezili útěku hráče.

\subsection{Větev Attacking}

Větev \texttt{Attacking} reprezentuje chování NPC v situaci, kdy je detekován nepřítel.
Je aktivována na základě hodnoty klíče \texttt{State} v \textit{blackboardu}, který je nastaven na \texttt{Attacking}
v rámci zpracování percepčních podnětů.

V této větvi NPC využívá klíč \texttt{AttackTarget}, který obsahuje referenci na aktuální cíl.
Na základě této hodnoty dochází k pohybu směrem k cíli a následnému provedení útoku.

Samotný útok je realizován pomocí úlohy \texttt{BTT\_Attack}, která spouští odpovídající animační montáž (viz \ref{fig:notify_timeline}) a řídí průběh útočné akce.

Chování v této větvi má nejvyšší prioritu, což znamená, že při detekci nepřítele dochází k okamžitému přepnutí z ostatních stavů do útočného režimu.

\subsection{Větev Investigating}

Větev \texttt{Investigating} reprezentuje chování NPC při vyhodnocení podezřelého podnětu,
typicky na základě sluchového vjemu.

Aktivace této větve je řízena hodnotou klíče \texttt{State} v \textit{blackboard},
který je nastaven na \texttt{Investigating} ve funkci \texttt{HandleSensedSound}.
Současně je do klíče \texttt{PointOfInterest} uložena poloha, odkud byl podnět zaznamenán.

V rámci této větve se NPC přesouvá k uložené pozici a provádí její prohledání.
Konkrétní průběh tohoto chování je řízen behaviorálním stromem na základě hodnoty klíče \texttt{PointOfInterest}.

Pokud během vyšetřování dojde k detekci nepřítele, dochází ke změně stavu na \texttt{Attacking}, což má vyšší prioritu a způsobí opuštění této větve.
Pokud se naopak ukáže, že prozkoumávané místo není ničím důležité, tak dojde ke změně do stavu \texttt{Passive}.

\subsection{Větev Passive}

Větev \texttt{Passive} reprezentuje výchozí chování NPC v situaci, kdy není detekován žádný významný podnět.
Slouží jako základní režim, ve kterém NPC vykonává hlídkování.

V rámci této větve je využita služba \texttt{BTS\_EnsurePatrolRoute}, která zajišťuje, že má NPC přiřazenou platnou hlídkovou trasu.
V případě, že žádná trasa není nastavena, služba vyžádá její přidělení prostřednictvím systému \texttt{BP\_RoutingManager}.

Samotný pohyb po trase je realizován úlohou \texttt{BTT\_MoveAlongPatrolRoute}, která postupně naviguje NPC mezi jednotlivými body hlídkové trasy.

Taktéž v této větvi probíhají interakce \texttt{BTT\_InteractWithActor}, přemisťování předmětů \texttt{BTT\_PlaceItem}, komunikace 
s jiným NPC \texttt{BTT\_InteractWithOtherNPC} a dělání rozhodnutí o změnách hlídkových tras. 

Pokud NPC nemá k dispozici žádnou trasu, setrvává na místě do doby, než je jí trasa přidělena.
Tato větev tak zajišťuje kontinuální chování NPC i v případě absence dalších podnětů z okolí.

\section{Implementace úloh a služeb}
\label{sec:implementace-tasky}

Tato část popisuje vybrané úlohy a služby behaviorálního stromu, které přímo propojují rozhodovací logiku stromu s konkrétními akcemi NPC ve světě hry. 
Úlohy typu \textit{Task} slouží k provedení jednorázové nebo déle trvající akce, zatímco služby typu \textit{Service} průběžně aktualizují data potřebná pro vykonávání dané větve stromu.

\subsection{BTT\_Attack}

Úloha \texttt{BTT\_Attack} realizuje spuštění útoku řízené NPC postavy. 
Při vykonání úlohy je nejprve získána reference na ovládanou postavu a ta je přetypována na \texttt{BP\_NPC}. 
Následně je na této postavě zavolána událost \texttt{Attack}, která je implementována v třídě \texttt{BP\_NPC} a zajišťuje přehrání útočné animační montáže.

Úloha není ukončena okamžitě po spuštění animace. 
Místo toho se naváže na událost \texttt{OnAttackEnd}, která je vyvolána až po dokončení útoku. 
Teprve poté je zavoláno \texttt{Finish Execute} s výsledkem \texttt{Success}. 
Tím je zajištěno, že behaviorální strom pokračuje až ve chvíli, kdy je útočná akce skutečně dokončena.

Součástí úlohy je také obsluha přerušení. 
Pokud je úloha předčasně ukončena, například kvůli přechodu do jiné větve stromu, je zastaven pohyb controlleru, zrušen aktuální focus 
a úloha je ukončena pomocí \texttt{Finish Abort}. 
Tento postup zabraňuje tomu, aby NPC po přerušení útoku pokračovalo v neaktuální akci.

\subsection{BTT\_MoveAlongPatrolRoute}

Úloha \texttt{BTT\_MoveAlongPatrolRoute} zajišťuje pohyb NPC po aktuálně přiřazené hlídkové trase. 
Při spuštění úloha načte z \textit{blackboardu} hodnotu klíče \texttt{CurrentPatrolRoute} a zkontroluje se, zda je klíč platný.

Z trasy je následně získána světová pozice aktuálního bodu pomocí funkce \texttt{GetSplinePointAsWorldPosition}. 
Na tuto pozici je NPC navigována pomocí uzlu \texttt{AI MoveTo}. 
Po úspěšném dosažení cíle je na trase zavolána funkce \texttt{IncrementPatrolRoute}, která posune interní ukazatel na další bod trasy. 
Poté je úloha ukončena pomocí \texttt{Finish Execute} s výsledkem \texttt{Success}.

Pokud se nepodaří získat platnou trasu, přetypování selže nebo pohyb k cílové pozici není úspěšný, úloha je ukončena neúspěšně. 
Součástí implementace je také reakce na přerušení úlohy, při které je pohyb NPC zastaven a úloha je ukončena pomocí \texttt{Finish Abort}.

\subsection{BTS\_EnsurePatrolRoute}

Služba \texttt{BTS\_EnsurePatrolRoute} zajišťuje, že má NPC během pasivního chování přiřazenou platnou hlídkovou trasu. 
Při aktivaci služby je z \textit{blackboardu} načten klíč \texttt{CurrentPatrolRoute}. 
Pokud je tato hodnota již nastavena, není nutné trasu měnit.

V případě, že NPC platnou trasu nemá, služba načte z \textit{blackboardu} referenci na systém \texttt{BP\_RoutingManager}. 
Po ověření platnosti této reference je řízená postava přetypována na \texttt{BP\_NPC} a je zavolána funkce \texttt{ChangeRoute}. 
Výsledkem této funkce je nová hlídková trasa, která je následně uložena zpět do klíče \texttt{CurrentPatrolRoute}.

Tato služba odděluje samotné přidělování tras od úlohy, která po trase pouze vykonává pohyb. 
Díky tomu může být logika výběru trasy spravována centralizovaně v \texttt{BP\_RoutingManager}, zatímco behaviorální strom pracuje pouze 
s aktuálně přiřazenou trasou.

\section{Implementace hlídkových tras}
\label{sec:implementace-routing}

\subsection{BP\_PatrolRoute}

Třída \texttt{BP\_PatrolRoute} reprezentuje hlídkovou trasu pomocí komponenty typu \texttt{SplineComponent}, která definuje plynulou křivku 
ve~světovém prostoru herní úrovně. 
Jednotlivé hlídkové body nejsou explicitně reprezentovány samostatnými komponentami, ale odpovídají bodům spline křivky.

Pro práci s~trasou třída poskytuje několik metod. 
Metoda \texttt{GetSplinePointAsWorldPosition} vrací světovou polohu bodu spline se zadaným indexem.
Dále je implementována metoda \texttt{IncrementPatrolRoute}, která aktualizuje index aktuálního bodu na základě směru pohybu a typu trasy.

Trasa může být buď cyklická, nebo otevřená. 
V případě cyklické trasy je index inkrementován o modulo počtu bodů křivky.
U otevřené trasy dochází při dosažení koncových bodů ke změně směru průchodu (tzv. ``ping-pong'' pohyb), což je realizováno přepínáním hodnoty 
proměnné \texttt{Direction}.

Pomocné metody \texttt{GetRouteProgress} a \texttt{GetRouteType}
umožňují získat relativní postup po trase a určit, zda se jedná
o cyklickou či necyklickou trasu.

\subsection{BP\_RoutingManager}

Třída \texttt{BP\_RoutingManager} implementuje správu hlídkových
tras a jejich přidělování NPC postavám.

Při inicializaci v rámci události \texttt{BeginPlay} jsou vyhledány všechny dostupné 
hlídkové trasy v levelu a ty jsou uloženy do interních kolekcí.
Trasy jsou dále rozděleny do skupin na základě tagů aktérů (např. \texttt{STAGE2}, \texttt{STAGE3}), což umožňuje řídit, které trasy jsou dostupné 
v různých fázích hry.

Obsazenost tras není reprezentována přímo v objektech tras, ale je evidována v mapě \texttt{Occupants}, která mapuje instance NPC na jim přidělené trasy.

Metoda \texttt{AcquireRoute(Requester)} vrací trasu pro daného žadatele. 
Pokud již má NPC trasu přidělenou, je vrácena tato existující trasa. 
V opačném případě dojde k náhodnému promíchání seznamu tras a následně k výběru první neobsazené trasy.
Vybraná trasa je zapsána do mapy \texttt{Occupants}.

Metoda \texttt{IsRouteOccupied(Route)} ověřuje, zda je daná trasa aktuálně obsazena některou NPC postavou.

Metoda \texttt{ReleaseRoute(Requester)} odstraní záznam o přiřazení trasy danému NPC a tím ji uvolní pro další použití.

Metoda \texttt{ChangeRoute(Requester)} nejprve uvolní aktuální trasu a následně přidělí novou trasu pomocí \texttt{AcquireRoute}.
Specifické varianty \texttt{ChangeRoute\_Stage2} a \texttt{ChangeRoute\_Stage3} vybírají trasu pouze z příslušné podmnožiny tras definované tagy.

Pokud není žádná volná trasa k dispozici, metoda vrací hodnotu \texttt{None}, přičemž NPC může o přiřazení požádat znovu v dalším průchodu stromem.

\section{Sdílení informací o předmětech}
\label{sec:sharing_information}

Komunikace mezi NPC není v implementaci řešena přímým předáváním informací o detekovaných hrozbách. 
Namísto toho je zaměřena na sdílení znalostí o herních předmětech, především o jejich výchozích pozicích a případném nesprávném umístění.

Základním prvkem tohoto systému je třída \texttt{BP\_GameStateBase}, která udržuje společný seznam předmětů ve scéně. Každý předmět odvozený
od třídy \texttt{BP\_BaseItem} se při spuštění zaregistruje do tohoto seznamu pomocí funkce \texttt{AddToItemList}. 
Současně je uložena jeho výchozí poloha prostřednictvím funkce \texttt{SetItemHomePosition}.
Tím vzniká sdílený zdroj informací, ke kterému mohou NPC při svém chování přistupovat.

Pro zjištění předmětu, který je potřeba uklidit, slouží funkce \texttt{RequestItemCleanup}. 
Ta prochází registrované předměty a hledá takový, který je označen jako nesprávně umístěný. 
Pokud je takový předmět nalezen, vrací referenci na něj společně s jeho aktuální polohou.

Vlastní výměna informací mezi NPC je realizována v úloze \texttt{BTT\_InteractWithOtherNPC}. 
Tato úloha vyhledá ostatní instance \texttt{BP\_NPC} ve scéně a vybere jinou postavu než tu, která úlohu právě vykonává. 
Následně pracuje se seznamy známých předmětů jednotlivých NPC a porovnává informace o jejich umístění.

Pokud NPC získá od jiné postavy relevantní informaci o předmětu, může tuto informaci využít pro aktualizaci vlastních hodnot v \textit{blackboardu}. 
V implementaci se tímto způsobem pracuje například s informací o ztraceném předmětu a s klíči souvisejícími s hledáním boudy. 
Díky tomu mohou NPC nepřímo sdílet informace o stavu herního světa, aniž by mezi nimi existoval samostatný systém zpráv nebo událostí.

\section{Implementační problémy a rozhodnutí}
\label{sec:problems_and_solutions}
Původní návrh systému komunikace mezi NPC počítal s individuální komponentou \texttt{BPC\_MemorySystem} připojenou ke každé NPC postavě. 
Tato komponenta měla uchovávat záznamy o objektech, se kterými postava přišla do kontaktu, a umožňovat přímou výměnu těchto záznamů mezi NPC.

Při implementaci se však ukázalo, že synchronizace paměťových záznamů mezi instancemi komponent napříč různými postavami nebyla spolehlivá. 
Docházelo k nekonzistentním stavům, kdy si NPC předávaly neaktuální nebo neplatné reference.

Problém byl vyřešen přechodem na centralizovaný přístup prostřednictvím třídy \texttt{BP\_GameStateBase}. 
Místo individuálních paměťových komponent byl zaveden společný registr objektů dostupný všem NPC. 
Každý objekt si při inicializaci zaregistruje výchozí polohu, ke které NPC přistupují při rozhodování. 
Přímá výměna paměťových záznamů byla nahrazena porovnáním lokálního seznamu \texttt{ItemsKnownList} s tímto centrálním registrem.



